\documentclass[11pt]{article}
\usepackage[margin=1in]{geometry}
\usepackage{amsmath,amssymb,amsthm,mathtools}
\usepackage[T1]{fontenc}
\usepackage{lmodern}
\usepackage{microtype}
\usepackage[hidelinks]{hyperref}
\usepackage{enumitem}
\newtheorem{theorem}{Theorem}
\newtheorem{lemma}{Lemma}
\newcommand{\N}{\mathbb N}
\newcommand{\R}{\mathbb R}
\newcommand{\Z}{\mathbb Z}
\newcommand{\Linfty}{L^{\infty}(\Omega,\mu)}
\newcommand{\code}[1]{\texttt{\small #1}}
\title{Quadratic ergodic averages need not converge in $L^\infty$\\
\large Disproof of conjecture 00000001925 as written}
\author{C0ldSmi1e}
\date{October 12, 2026}
\begin{document}
\maketitle
\begin{abstract}
We construct one invertible probability-preserving transformation and one bounded
measurable observable for which the quadratic averages satisfy
$\|A_{2N}f-A_Nf\|_\infty\geq 1/2$ for every $N\geq1$.
Thus these averages have no $L^\infty$ limit, so no rate tending to zero can hold
for the unqualified class in the conjecture. The example is nonergodic;
a different statement requiring the transformation to be ergodic is not settled here.
\end{abstract}

\section{Original claim and precise scope}
The English original reads:
\begin{quote}
Conjecture: The polynomial ergodic averages converge in L\textsuperscript{$\infty$}
norm at a polylogarithmic multiple-logarithm rate (a Bernstein-type bound for p-averages).
\end{quote}
Neither this text nor its Chinese version specifies a rate formula or assumes
that the transformation itself is ergodic. We use the usual meaning of
\emph{polynomial ergodic averages}: on a probability space, for a
measure-preserving transformation $T$, a bounded measurable real function $f$,
and an integer-valued polynomial $p$, the averages are
\[
 A_Nf(x)=\frac1N\sum_{k=0}^{N-1} f(T^{p(k)}x),\qquad N\geq1.
\]
We choose $p(k)=k^2$, a genuine degree-two polynomial. ``Ergodic averages''
names the averaging operation and does not impose ergodicity of $T$.
The theorem below disproves qualitative norm convergence under these stated
standard assumptions. It does not assign an invented meaning to the unspecified
rate or prove anything about a stronger hypothesis of ergodicity.

\section{One fixed probability-preserving system}
For $n\in\N=\{0,1,2,\ldots\}$ let
\[
 q_n=4(n+1)^2+1,\qquad
 \Omega=\bigsqcup_{n\geq0}\bigl(\{n\}\times\Z/q_n\Z\bigr),\qquad
 w_n=2^{-(n+1)}.
\]
Equip this countable set with its full power-set sigma-algebra and put
\[
 \mu\{(n,j)\}=\frac{w_n}{q_n}.
\]
All atoms have positive measure, and summing first within a component gives
$\mu(\Omega)=\sum_{n\geq0}w_n=1$.
Define
\[
 T(n,j)=(n,j+1),\qquad
 f(n,j)=\mathbf1_{\{[j]_{q_n}\geq(n+1)^2\}},
\]
where $[j]_{q_n}\in\{0,\ldots,q_n-1\}$ is the canonical residue.
Both maps are measurable. The inverse of $T$ is $(n,j)\mapsto(n,j-1)$.
Rotation permutes each finite component, whose atoms have equal mass, so
$T$ preserves $\mu$. Also $0\leq f\leq1$, so $f\in\Linfty$.
All this data is fixed once and for all, independently of $N$.

\section{A uniform gap at every positive length}
\begin{lemma}
For every $N\geq1$, set $x_N=(N-1,0)$. Then
$A_Nf(x_N)=0$ and $A_{2N}f(x_N)=1/2$.
\end{lemma}
\begin{proof}
The component of $x_N$ has length $4N^2+1$. For $0\leq k<2N$,
$k^2<4N^2+1$, so $T^{k^2}x_N=(N-1,k^2)$ has not wrapped around the cycle.
The threshold defining $f$ on this component is $N^2$. Thus
\[
 f(T^{k^2}x_N)=
 \begin{cases}
 0,&0\leq k<N,\\
 1,&N\leq k<2N.
 \end{cases}
\]
The first average contains $N$ zeros. The doubled average contains $N$ zeros
and $N$ ones, and its denominator is $2N$.
\end{proof}

\begin{theorem}
In the genuine $L^\infty$ quotient space,
\[
 \|A_{2N}f-A_Nf\|_\infty\geq\frac12\quad(N\geq1).
\]
Consequently $(A_Nf)$ is not Cauchy in $L^\infty$ and does not converge to any
$L^\infty$ function.
\end{theorem}
\begin{proof}
At $x_N$ the absolute difference equals $1/2$, and
\[
 \mu\{x_N\}=\frac{2^{-N}}{4N^2+1}>0.
\]
Therefore this value cannot be removed by discarding a null set: the essential
supremum of the difference is at least $1/2$. Although the witness atom depends
on $N$ and its mass tends to zero, the space, measure, transformation, and
observable do not depend on $N$.
If the averages were Cauchy, beyond some index all pairwise distances would
be strictly less than $1/2$. Taking $N$ beyond that index contradicts the displayed
gap between $N$ and $2N$. Every convergent sequence in a normed space is Cauchy,
so convergence to any possible limit is excluded.
\end{proof}

The invariant component $C_0=\{(0,j):j\in\Z/5\Z\}$ has measure $w_0=1/2$;
this verifies the advertised nonergodicity. No claim is made about convergence
for an amended problem restricted to ergodic transformations. The example
already refutes norm convergence for the unqualified probability-preserving
class, hence also any proposed norm-convergence rate on that class.

\clearpage
\section{Formal correspondence and verification}
The accompanying project uses Lean 4.19.0 and Mathlib v4.19.0, with exact
transitive dependency revisions in \code{lake-manifest.json}. Its core definitions
are actual Mathlib objects, not surrogate numerical sequences:
\begin{itemize}[leftmargin=1.5em,itemsep=3pt]
\item \code{Dynamics.lean} defines the dependent sum of finite cyclic groups,
the invertible rotation, the observable, and normalized sums of its actual
square-time iterates. It proves the two values in Lemma 1 and pointwise bounds.
\item \code{Measure.lean} defines a normalized probability mass function using
geometric component weights, converts it to a measure, and proves positive
singleton masses and preservation of that same measure by the rotation.
\item \code{Nonconvergence.lean} proves essential boundedness and uses
\code{MemLp.toLp} to put the averages in Mathlib's \code{Lp} space with exponent
$\infty$. An essential-supremum estimate gives the quotient norm gap, followed by
failure of the Cauchy condition and exclusion of every limit.
\item \code{Counterexample.lean} defines the main theorem, \code{counterexample},
which bundles the full statement and the explicit nonergodicity clause.
\code{Tests.lean} contains 15 kernel-checked integration examples. The library
root imports the tests and thereby includes every authored proof in the default build.
\end{itemize}
On this countable space every singleton has positive mass, so an almost-everywhere
predicate holds exactly when it holds at every point. This is the formal bridge
from the witness calculation to the essential supremum. The proof covers every
positive integer length; no finite numerical sample is used as evidence for a
universal claim. No auxiliary numerical computation is required.

\paragraph{Recorded local validation.}
The default Lake build and fresh replay of all six authored modules and the Lake
configuration passed, with warnings treated as errors. All five designated
capstone theorems have only the axiom dependencies \code{propext},
\code{Classical.choice}, and \code{Quot.sound}. There are no authored admissions,
custom axioms, or native-decision shortcuts in the mathematical project.

A declaration audit recorded 96 originating declarations, a dependency graph of
20,672 declarations, and 4,063 resolved imported modules. All 74 kernel-safe
originating declarations passed the permitted-axiom and safe-dependency policy.
The remaining 22 entries are generated runtime helpers, including two unsafe
lambda-lifting axiom entries created by Lean's compiler. Their provenance was
checked against the exact Lean 4.19.0 compiler source, and none is reachable from
any of the 74 safe declaration closures. They are not assumptions of the proof.
Six control fixtures checked that the audit accepts an ordinary proof and rejects
custom/transitive axioms, an admitted proof, native decision, and unsafe code.

A separate local semantic review checked the exact bilingual original, the full
frozen source, and the complete mathematical report. The PDF was compiled and
visually checked. The submission includes the verification tools, source hashes,
review records, and compressed raw audit evidence. All mathematical identities
are proved symbolically; no numerical helper program is needed for this disproof.

\paragraph{Trust boundary.}
All nine dependency revisions and their tracked sources were checked against the
pins; cached imports were matched to the supplied pinned stock. Full Mathlib was
not rebuilt, and the supplied Lean runtime was not bootstrapped. The dependency
inspector is an audit tool, not an independent kernel rechecker; source compilation
supplies kernel checking. These are local validation results, not a claim of
maintainer acceptance.
\end{document}
